\documentclass[10pt,a4paper]{amsart}
\usepackage[margin=19mm]{geometry}
\usepackage{amsmath,amssymb,mathtools}
\usepackage[scr=rsfs,cal=euler]{mathalfa}
\usepackage{xfrac}
\usepackage[unicode,hidelinks,bookmarks=false]{hyperref}
\newcommand{\ie}{i.e.\ }
\newcommand{\cf}{cf.\ }
\newcommand{\resp}{respectively\ }
\newcommand{\Subsection}{\subsection}
\providecommand{\nobreakdash}{\nobreak\hbox{-}}
\setlength{\emergencystretch}{2em}
\multlinegap0pt
\usepackage{bm}
\usepackage{bbm}
\usepackage{dsfont}
\usepackage{xspace}
\usepackage{cancel}
\usepackage{mathtools}
\usepackage[shortlabels]{enumitem}
\usepackage{amsthm}
\makeatletter
\@ifundefined{theorem}{\newtheorem{theorem}{Theorem}}{}
\@ifundefined{corollary}{\newtheorem{corollary}[theorem]{Corollary}}{}
\@ifundefined{lemma}{\newtheorem{lemma}[theorem]{Lemma}}{}
\makeatother
\newcommand{\calD}{\mathcal{D}}
\renewcommand{\Pr}{\mathbb{P}}
\title[Coloring small locally sparse degenerate graphs and related ]{Coloring small locally sparse degenerate graphs and related problems}
\author{Domagoj Brada\v{c}, Jacob Fox, Raphael Steiner, Benny Sudakov, Shengtong Zhang}
\date{Source: arXiv:2601.15245v1 (CC BY 4.0)}
\begin{document}
\maketitle

\noindent\emph{Source and licence.} Coloring small locally sparse degenerate graphs and related problems. Version arXiv:2601.15245v1 (CC BY 4.0), distributed under the Creative Commons Attribution 4.0 International licence, as recorded in the arXiv licence metadata for this exact version and shown on the abstract page https://arxiv.org/abs/2601.15245v1. Full rights notice: licenses/arxiv-2601.15245-CC-BY-4.0.txt.\par\smallskip

\noindent\textit{Editorial scope.} Counted unit: the Corollary whose source label is \texttt{cor:sparseleft} in the article (source0.tex lines 340--348, source label \texttt{cor:sparseleft}), delivered together with the article's own complete original proof of it (source0.tex lines 349--373, one \texttt{proof} environment of 25 lines). No mathematics is written, repaired or re-derived: statement and proof are the article's own text line for line. Required context retained uncounted: lem:main (lines 243--252); prop:indepofwhatcomesafter (lines 246--251). Every definition, display and label of those lines is retained unchanged. Omitted from this package and not redistributed: 1--242, 252--339, 374--801. Nothing inside a retained excerpt is omitted or altered: every retained body is delivered exactly as the article's lines are written, so no omission receipt of any kind accompanies this package. Citations: the retained excerpts carry no citation command at all, so no bibliography entry of the article is transcribed and no reference is redistributed. Reference adaptation: none was needed and none was made; every reference inside the delivered unit resolves inside it, so no unresolved reference or citation is produced. Printed slips kept literal: no printed word, symbol, hypothesis or number of the counted statement or of its proof was altered. Layout and class: the article's own class is \texttt{amsart} with packages including graphicx, geometry, hyperref, amsthm, complexity, thm-restate, amssymb, amsmath, mathtools, cleveref, comment, setspace, xcolor, enumitem; the wrapper is the lane's pinned \texttt{amsart} editorial wrapper at 10pt on A4, which restates the theorem-like environments the retained text needs and adds the article's own macro definitions that the retained text needs (\texttt{\textbackslash Pr}, \texttt{\textbackslash calD}), with extra wrapper packages (bm, bbm, dsfont, xspace, cancel, mathtools, enumitem). The delivered numbering is the unit's own. Assets: no retained span contains a figure environment, a graphics-inclusion command, a PDF-page inclusion, a table or a TikZ picture; no image file is copied into this package, the package declares no asset dependency and no image file is redistributed with it. Licence: version arXiv:2601.15245v1 of this article is distributed under the Creative Commons Attribution 4.0 International licence, as recorded in the arXiv licence metadata for this exact version and shown on the abstract page https://arxiv.org/abs/2601.15245v1. A full rights notice is delivered at licenses/arxiv-2601.15245-CC-BY-4.0.txt. Characters: every retained excerpt is pure ASCII, so the delivered engine, XeLaTeX with its default Latin Modern text font, reports no missing character.
\begin{lemma} \label{lem:main}
    Let $d\in \mathbb{R}_{\ge 1}$, $0 < \alpha < \frac{\log d}{4d}$ and $G$ be a $d$-degenerate triangle-free graph. Let $v_1,\ldots,v_n$ be a linear ordering witnessing that $G$ is $d$-degenerate, i.e., such that $|N_L(v_i)|\le d$ for all $i\in [n]$. Then there exists a probability distribution $\mathcal{D}$ on the independent sets of $G$ with the following properties.

    \begin{enumerate}[label=\alph*)]
        \item \label{prop:large-prob} For every $i \in [n], \, \Pr_{I \sim \calD}[v_i \in I] \ge \alpha / 4.$
        \item \label{prop:prob-not-too-big} For every $i \in [n]$ and every independent set $S \subseteq \{v_1, \dots, v_{i-1}\},$ it holds that
        \[ \Pr_{I \sim \calD}[v_i \in I \, \vert \, I \cap \{v_1, \dots, v_{i-1}\} = S] \le \alpha \cdot e^{2 \alpha d}. \]
        \item \label{prop:indepofwhatcomesafter} For every $i\in [n]$, the distribution of the set $I\cap \{v_1,\ldots,v_i\}$ where $I$ is sampled from $\mathcal{D}$ depends only on the ordered induced subgraph $G[\{v_1,\ldots,v_{i}\}]$ of $G$.
    \end{enumerate}
\end{lemma}

\begin{corollary} \label{cor:sparseleft}
    Let $d\in \mathbb{R}_{\ge 1}$ and let $G$ be a $d$-degenerate graph. Let $v_1,\ldots,v_n$ be a linear ordering witnessing that $G$ is $d$-degenerate and let $f\in (2,d^2)$ be such that $|E(G[N_L(v_i)])|\le d^2/f$ for every $i \in [n]$. Finally, let $0<\alpha<\frac{\log f}{8\sqrt{f}}$ be a parameter. Then there exists a probability distribution $\mathcal{D}$ on the independent sets of $G$ satisfying the following.

    \begin{enumerate}[label=\alph*)]
        \item \label{prop:large-prob2} For every $i \in [n], \, \Pr_{I \sim \calD}[v_i \in I] \ge \frac{\alpha\sqrt{f}}{32 d}.$
        \item \label{prop:prob-not-too-big2} For every $i \in [n]$ and every independent set $S \subseteq \{v_1, \dots, v_{i-1}\},$ it holds that
        \[ \Pr_{I \sim \calD}[v_i \in I \, \vert \, I \cap \{v_1, \dots, v_{i-1}\} = S] \le \frac{\alpha\sqrt{f} \cdot e^{2 \alpha \sqrt{f}}}{2d}. \]
    \end{enumerate}
    \end{corollary}

    \begin{proof}
    Set $p:=\frac{\sqrt{f}}{2d}\in (0,1)$ and let $U\subseteq V(G)$ be a random subset of vertices of $G$ created according to the following process: First, we create a random subset $V\subseteq V(G)$ in which every vertex of $G$ is included independently with probability exactly $p$. Then, a vertex $v_i$ is put into $U$ if and only if $v_i\in V$, $|N_L(v_i)\cap V|\le \sqrt{f}$ and $N_L(v_i)\cap V$ is an independent set in $G$. 

    Let $i\in [n]$ be arbitrary and let us estimate the probability that $v_i\in U$. Note that by Markov's inequality, we have that
    $$\mathbb{P}[|N_L(v_i)\cap V|>\sqrt{f}]<\frac{p|N_L(v_i)|}{\sqrt{f}}\le \frac{
pd}{\sqrt{f}}=\frac{1}{2}.$$
Similarly, we have
$$\mathbb{P}[N_L(v_i)\cap V \text{ not independent}]\le \mathbb{E}[|E(G[N_L(v_i)\cap V])|]=p^2|E(G[N_L(v_i)])|\le p^2\frac{d^2}{f}=\frac{1}{4}.$$
From this, we immediately conclude that $\mathbb{P}[v_i\in U]\ge \mathbb{P}[v_i\in V](1-\frac{1}{2}-\frac{1}{4})=\frac{p}{4}$ for every $i\in [n]$. 
    
    Note that by definition of the process to generate $U$, we will always have that the induced subgraph $G[U]$ is triangle-free and $d':=\sqrt{f}$-degenerate (the degeneracy is witnessed by the linear ordering obtained from $v_1,\ldots,v_n$ by omitting the vertices in $V(G)\setminus U$). Since $0<\alpha<\frac{\log f}{8\sqrt{f}}= \frac{\log d'}{4d'}$, we can apply Lemma~\ref{lem:main} to $G[U]$. Let $\mathcal{D}_U$ be the probability distribution on independent sets in $G[U]$ guaranteed by the lemma. Fixing $U$, from the properties guaranteed by the lemma, we then have
    $$\mathbb{P}_{I\sim \mathcal{D}_U}[v_i\in I]\ge \frac{\alpha}{4}$$ for every $v_i\in U$, and for every independent set $S\subseteq \{v_1,\ldots,v_{i-1}\}\cap U$ in $G$, it holds that
    $$\mathbb{P}_{I\sim \mathcal{D}_U}[v_i\in I|I\cap \{v_1,\ldots,v_{i-1}\}=S]\le \alpha\cdot e^{2\alpha d'}=\alpha\cdot e^{2\alpha\sqrt{f}}.$$

    Finally, let us consider the distribution $\mathcal{D}$ on independent sets of $G$ described by first sampling $U$ randomly according to the above process and then sampling an independent set from $\mathcal{D}_U$. We claim that it satisfies the desired properties. Indeed, for every $i\in [n]$, we have that
    $$\mathbb{P}_{I\sim \mathcal{D}}[v_i\in I]=\sum_{A\subseteq V(G):v_i\in A}\mathbb{P}_{I\sim \mathcal{D}}[v_i\in I|U=A]\cdot\mathbb{P}[U=A]=\sum_{A\subseteq V(G):v_i\in A}\mathbb{P}_{I\sim \mathcal{D}_A}[v_i\in I]\cdot\mathbb{P}[U=A]$$
    $$\ge \frac{\alpha}{4}\sum_{A\subseteq V(G): v_i\in A}\mathbb{P}[U=A]=\frac{\alpha}{4}\mathbb{P}[v_i\in U]\ge \frac{\alpha}{4}\cdot \frac{p}{4}=\frac{\alpha\sqrt{f}}{32d},$$ as desired. Next, consider any fixed independent set $S\subseteq \{v_1,\ldots,v_{i-1}\}$ in $G$. 
    We then have
    $$\mathbb{P}_{I\sim \mathcal{D}}[v_i\in I|I\cap \{v_1,\ldots,v_{i-1}\}=S]$$ $$=\sum_{\substack{A\subseteq V(G):\\ v_i\in A, S\subseteq A}}\mathbb{P}_{I\sim \mathcal{D}}[v_i\in I|I\cap \{v_1,\ldots,v_{i-1}\}=S, U=A]\cdot\mathbb{P}[U=A|I\cap \{v_1,\ldots,v_{i-1}\}=S]$$
    $$=\sum_{\substack{A\subseteq V(G):\\ v_i\in A, S\subseteq A}}\mathbb{P}_{I\sim \mathcal{D}_A}[v_i\in I|I\cap \{v_1,\ldots,v_{i-1}\}=S]\cdot\mathbb{P}[U=A|I\cap \{v_1,\ldots,v_{i-1}\}=S]$$
    $$\le \alpha e^{2\alpha\sqrt{f}}\sum_{\substack{A\subseteq V(G):\\ v_i\in A, S\subseteq A}}\mathbb{P}[U=A|I\cap \{v_1,\ldots,v_{i-1}\}=S]$$
    $$=\alpha e^{2\alpha\sqrt{f}} \mathbb{P}[v_i\in U|I\cap \{v_1,\ldots,v_{i-1}\}=S]$$
    $$\le \alpha e^{2\alpha\sqrt{f}} \mathbb{P}[v_i\in V|I\cap \{v_1,\ldots,v_{i-1}\}=S].$$
    Finally, note that the events $v_i\in V$ and $I\cap \{v_1,\ldots,v_{i-1}\}=S$ are independent. This is because, by definition of $V$, the event $v_i\in V$ is independent from $V\cap \{v_1,\ldots,v_{i-1}\}$, which by definition fully determines the set $U\cap \{v_1,\ldots,v_{i-1}\}$. Finally, by Lemma~\ref{lem:main},~\ref{prop:indepofwhatcomesafter} the latter set fully determines the distribution of $I\cap \{v_1,\ldots,v_{i-1}\}$ when $I$ is sampled from $\mathcal{D}_U$. Hence, we find that $$\mathbb{P}_{I\sim \mathcal{D}}[v_i\in I|I\cap\{v_1,\ldots,v_{i-1}\}=S]\le \alpha e^{2\alpha\sqrt{f}}\mathbb{P}[v_i\in V]=\alpha e^{2\alpha\sqrt{f}}p=\frac{\alpha\sqrt{f} e^{2\alpha\sqrt{f}}}{2d},$$ as desired. This concludes the proof of the corollary.
    \end{proof}

\end{document}
